\section{Introduction \& the problem this replaces}\label{sec:intro}

QtRocket is an open-source model-rocket simulator built on a shared C++23 simulation core that
drives both a Qt6 Widgets GUI and a headless command-line REPL, with a standalone OpenGL viewer
sharing the same rocket model. This document specifies a replacement for the way that core
expresses \emph{where one part sits relative to another}. The replacement is structural rather
than cosmetic: it removes a coordinate convention that is correct only by accident and replaces it
with one that is correct by construction.

\subsection{How parts are linked today}

A rocket is a composite tree of \code{Part} objects (concrete types such as \code{ConicalNoseCone},
\code{BodyTube}, \code{FinSet}, and \code{HollowSphere} live under \srcf{model/parts/}). Each node
stores its direct children together with a single relative-position vector. The storage is, at
\src{model/parts/Part.h}{326},

\begin{minted}{cpp}
std::vector<std::tuple<std::shared_ptr<Part>, Vector3>> childParts;
\end{minted}

\noindent and the doc comment on that member states the contract precisely: the \cppinline|Vector3|
is ``the relative position of the child part's center of mass with respect to the center of mass of
this part.'' In other words, parts are linked \emph{center-of-mass to center-of-mass} (CM-to-CM).
That single offset is the only spatial relationship the model records; everything downstream --- the
composite mass properties consumed by the integrator and the geometry drawn by the viewer --- is
derived from it.

\classfig{class-current}{The current composite tree links children CM-to-CM: each edge carries a
single \texttt{Vector3} measured from parent CM to child CM, and both the simulator and the
visualizer reconstruct absolute position by accumulating those offsets.}{fig:class-current}{0.86}

\subsection{Why a CM-to-CM link is fragile}

The CM-to-CM link rests on a hidden assumption: that a part's center of mass is a stable, intuitive
landmark an author can reason about. For a longitudinally symmetric part it happens to be ---
the CM is at mid-length --- but \cref{fig:class-current} hides the fact that \emph{this is not true
in general}. A solid right cone has its centroid one quarter of its length from the base; a thin
conical shell has its centroid one third of its length from the base. QtRocket encodes exactly this
in \src{model/parts/ConicalNoseCone.cpp}{49}:

\begin{minted}{cpp}
// With the cone in [-L/2, +L/2] (base at +L/2), the centroid is hbar forward of the base:
// L/4 (solid), L/3 (shell) => z = +L/2 - hbar.
Vector3 ConicalNoseCone::coneCmOffset(double L, bool solid)
{
   const double hbar = solid ? (L / 4.0) : (L / 3.0);
   return Vector3{0.0, 0.0, L / 2.0 - hbar};
}
\end{minted}

\noindent so the cone's CM is displaced from its geometric center (mid-length) by $L/2 - \bar{h} =
L/4$ for a solid cone and $L/6$ for a shell. An author who wants a nose cone's base rim to meet a
body tube's fore rim must therefore not state that physical intent directly; instead they must
hand-compute the offset between two centers of mass that sit at different fractions of their
respective lengths. The link records a derived quantity (a CM displacement) where the author
actually knows a geometric fact (two rims touch). Every such offset is a manual calculation, and
every manual calculation is a place to be wrong.

\begin{keyinsight}
A CM-to-CM link is only as trustworthy as the assumption that CM coincides with the geometric
center. That assumption holds for every symmetric part and fails for exactly the part most rockets
begin with --- the nose cone.
\end{keyinsight}

\subsection{The silent divergence between simulator and visualizer}

The deeper failure is that the two consumers of the link do not even agree on what the offset
\emph{means}. The simulation core treats it as documented: in \src{model/parts/Part.cpp}{169},
\code{Part::computeCompositeAt} accumulates each child's contribution at \cppinline|pos + cc.cm| ---
the stored offset \emph{plus} the child subtree's own CM --- and then shifts every tensor to the
composite CM via the parallel-axis theorem. This is a strictly CM-to-CM computation.

The visualizer makes a different assumption entirely. In \srcf{visualizer/RocketMesh.cpp}, the
\code{RocketMesh::walk} recursion (the lambda at line~427) builds each primitive
\emph{geometrically centered} --- \code{buildCone} and \code{buildTube} place their mesh symmetrically
about the part's mid-length --- and then translates it by an accumulated \code{station}:

\begin{minted}{cpp}
const Vector3& offset = std::get<1>(childPair);   // the stored CM-to-CM Vector3
if(child)
{
   self(self, *child, station + offset);          // accumulated as geometric-center-to-center
}
\end{minted}

\noindent The walk's own comment names the discrepancy: ``\code{station} = accumulated CM-to-CM
offsets from the root down to this part.'' The renderer takes a value \emph{defined} as a CM-to-CM
displacement and \emph{applies} it as a geometric-center-to-geometric-center translation. For any
symmetric part the two interpretations coincide and nothing is visibly wrong. For a nose cone they
diverge by precisely the cone's $(\text{geometric center} - \text{CM})$ offset --- $L/4$ for a solid
cone --- and so the simulator and the viewer silently reconstruct two different rockets from the same
file.

\begin{hardfail}
This is not hypothetical. The fixture \src{tests/data/designs/xl75\_multi.qrd}{19} places a coupler
at \mintinline{xml}{<offset x="0" y="0" z="0.48999999999999999"/>} relative to its body. Read as a
CM-to-CM offset by the simulator and as a geometric translation by the viewer --- and with the
nose's own CM-vs-center displacement folded in on top --- the link produces a coupler that
interpenetrates the nose cone. The two consumers disagree, no check fires, and the design loads,
simulates, and renders without complaint. \fail
\end{hardfail}

\subsection{The fix, in one paragraph}

The remedy is to remove the divergence by construction. Inter-part relationships are stated as
physical \emph{intent} --- ``this rim seats against that rim,'' ``this tube nests inside that bore''
--- expressed as a pair of fractional axial stations plus a seat kind, never as an opaque offset. A
\textbf{single resolver} walks the ownership tree and produces an absolute, purely \emph{geometric}
placement for every part, and \emph{both} the simulator (\code{computeCompositeAt},
\code{accumulateAeroAt}) and the visualizer (\code{RocketMesh::walk}) consume that one placement.
Because there is now exactly one source of absolute position, the two consumers cannot disagree.
Center of mass ceases to be an input and becomes a purely \emph{derived} consequence of geometry:
the resolver places parts by their edges, and the existing two-pass mass-weighted CM and
parallel-axis machinery is re-driven from those geometric stations rather than from author-supplied
CM offsets. No force-path or integrator change is required, and a path to full six-degree-of-freedom
placement is shaped in at zero schema cost.

\subsection{Provenance of this specification}

This specification is not a single author's first draft. It is the adjudicated synthesis of
\textbf{five independent design sketches} --- feature-mating, datum-pair, declarative
constraint-solver, local-frame scene graph, and joint-as-entity graph --- judged on three axes
(correctness and robustness; simplicity and migration cost; extensibility and ergonomics) and then
hardened against an adversarial critique. The judges converged on the datum/station-pair chassis,
upgraded it with the extent-based overlap rigor of the constraint-solver sketch, the \code{Pose} and
snap-operator mechanics of the scene-graph sketch, and the direction-constrained nesting of the joint
sketch. The rejected alternatives, and \emph{why} each was rejected, are recorded throughout this
document rather than discarded, because the rejections are themselves part of the argument for the
chosen design.

\subsection{Document roadmap}

The remainder of this document develops the design in the order an implementer would build it.
\Cref{sec:philosophy} states the three non-negotiable design rules and settles the longitudinal
convention (\zfwd, toward the nose tip), under which every part --- the cone included --- shares one
local frame, correcting a reversed-frame defect in the current cone rather than adapting around it.
\Cref{sec:datamodel} gives the exact C++ value types
(\code{StationLink}, \code{Station}, \code{SeatKind}) and the in-slot storage swap that replaces the
\cppinline|(Part, Vector3)| tuple. \Cref{sec:resolver} specifies the single resolver, how CM becomes
derived, and the caching split that keeps placement out of the per-step mass gate.
\Cref{sec:authoring} covers the authoring API and the zero-config default; \cref{sec:diagnostics}
the three-layer overlap diagnostics that turn the silent interpenetration above into a hard, located
failure. \Cref{sec:serialization} updates the \code{.qrd} schema, \cref{sec:migration} lays out the
physics-invariant phased migration, and \cref{sec:sixdof} shows the six-degree-of-freedom path the
design is shaped for. \Cref{sec:worked} works the \code{xl75\_multi} fixture end to end;
\cref{sec:deferred} records what was deliberately deferred or not adopted; and \cref{sec:files},
\cref{sec:repro}, and \cref{sec:glossary} enumerate the touched files, the reproduction steps, and
the glossary.
